
\subsection{Latent variable model}

In cases where $Y^{obs}=0$, we estimate the probability that $Y=1$ using a latent variable model. 

\subsubsection{Parameterization for unobserved outcoome.}

For the treated group, we parameterize the probability of being alive and stably housed at follow-up as a function of baseline covariates, treatment initiation time, and program enrollment status by marginalizing across the strata that include stable housing (i.e.\ $C \in \{3,4,5\}$) as well as over SSVF enrollment status at follow-up:

\begin{align*}
    P[&Y=1 \mid A=1, T^A=s, \boldsymbol{L}, \boldsymbol{X}(0), H(0)] \\
    P[&D(\tau)=0, H(\tau)=3 \mid A=1, T^A=s, \boldsymbol{L}, \boldsymbol{X}(0), H(0)] \\
     =& P[D(\tau)=0, H(\tau)=3 \mid A=1, T^A=s, \boldsymbol{L}, \boldsymbol{X}(0), H(0)] \\
     =& \sum_{c=3}^5 P[H(\tau)=3 \mid A=1, T^A=s, \boldsymbol{L}, \boldsymbol{X}(0), H(0), C=c]P[C=c \mid A=1, T^A=s, \boldsymbol{L}, \boldsymbol{X}(0), H(0)] \\
     =& \sum_{c=3}^5 \sum_{z} P[H(\tau)=3 \mid A=1, T^A=s, Z=z, \boldsymbol{L}, \boldsymbol{X}(0), H(0), C=c]\\
     & \quad \quad \quad \times P[C=c \mid A=1, T^A=s, Z=z \boldsymbol{L}, \boldsymbol{X}(0), H(0)] \\
     &\quad \quad \quad \times P[Z=z \mid A=1, T^A=s, \boldsymbol{L}, \boldsymbol{X}(0), H(0)]
\end{align*}

For the untreated group, we do the same thing but without marginalizing over SSVF status:
\begin{align*}
    P[&Y=1 \mid A=0, T^A=s, \boldsymbol{L}, \boldsymbol{X}(0), H(0)] \\
    P[&D(\tau)=0, H(\tau)=3 \mid A=0, T^A=s, \boldsymbol{L}, \boldsymbol{X}(0), H(0)] \\
     =& P[D(\tau)=0, H(\tau)=3 \mid A=0, T^A=s, \boldsymbol{L}, \boldsymbol{X}(0), H(0)] \\
     =& \sum_{c=3}^5 P[H(\tau)=3 \mid A=0, T^A=s, \boldsymbol{L}, \boldsymbol{X}(0), H(0), C=c]P[C=c \mid A=0, T^A=s, \boldsymbol{L}, \boldsymbol{X}(0), H(0)] \\
     =& \sum_{c=3}^5  P[H(\tau)=3 \mid A=0, T^A=s \boldsymbol{L}, \boldsymbol{X}(0), H(0), C=c]\\
     & \quad \quad \quad \times P[C=c \mid A=0, T^A=s, Z=z \boldsymbol{L}, \boldsymbol{X}(0), H(0)]
\end{align*}


For both groups, we assume that $H(\tau)$ is conditionally independent of baseline covariates given stratum membership, treatment status, and a subset of covariates $\boldsymbol{B} \subseteq [\boldsymbol{L}, \boldsymbol{X}(0)]$.

\begin{align*}
    P[&H(\tau)=h \mid A=a, T^A=s, Z=z, \boldsymbol{L}, \boldsymbol{X}(0), H(0), C=c] \\
    &= P[H(\tau)=h \mid A=a, T^A=s, \boldsymbol{B}, H(0), C=c]  
\end{align*}



\subsubsection{Parameterization for visit process.}

At baseline, we assume the following model for the visit process during an initial time window of 15 days after NCCHV contact:
\begin{align*}
    &\frac{P[V(t)=v \mid \boldsymbol{L}, \boldsymbol{X}(0), H(0)]}{P[V(t)=0 \mid \boldsymbol{L}, \boldsymbol{X}(0), H(0)]} \\
    &= \delta_0^v + \boldsymbol{\delta}_1^v \cdot \boldsymbol{L} + \boldsymbol{\delta}_2^v \cdot \boldsymbol{X}(0) + \delta_3^v \cdot H(0) + \boldsymbol{\delta}_4^v f(H(0), \boldsymbol{X}(0), \boldsymbol{L}) \quad \text{for } t \in [0, 15] 
\end{align*}
where $f(H(0), \boldsymbol{X}(0), \boldsymbol{L})$ is a vector of interaction terms between baseline housing status, baseline covariates, and baseline program enrollment. We assume that the visit process during the pre-outcome observation window $[t^*, \tau]$ follows the same model but with additional parameters to allow for changes in the visit process over time and differential changes across strata:

\begin{align*}
    &\frac{P[V(t)=v \mid \boldsymbol{L}, \boldsymbol{X}(0), H(t)]}{P[V(t)=0 \mid \boldsymbol{L}, \boldsymbol{X}(0), H(t)]} \\
    &= \delta_0^v + \boldsymbol{\delta}_1^v \cdot \boldsymbol{L} + \boldsymbol{\delta}_2^v \cdot \boldsymbol{X}(0) + \delta_3^v \cdot H(t) + \boldsymbol{\delta}_4^v f(H(t), \boldsymbol{X}(0), \boldsymbol{L}) + \delta_5^v C +  \delta_6^v H(t) + \delta_7^v C \times H(t) +  \quad \text{for } t \in [t^*, \tau]
\end{align*}


\subsubsection{Parameterization for manifest variables.}
Within each stratum and each visit setting, we paramterize each of the manifest variables as a function of baseline covariates, treatment status, and program enrollment status. Similarly to the visit process, we assume that the measurement model during the followup observation window $[t^*, \tau]$ follows the same model as at baseline but with additional parameters to allow for changes in the measurement process over time and differential changes across strata.

For $v=1, 2$, we parameterize ICD-10 codes  $M_1^v(t)$ 


Within each stratum $C=c$ for $c\in\{3,4,5\}$, we construct a series of estimating equations that relate the manifest variables observed in $[t^*,\tau]$ to the latent true housing status at $\tau$. We then solve these equations to estimate $P[H(\tau)=3\mid T^A=s, C=c]$ for treated patients and $P[H(\tau)=3\mid A=0, C=c]$ for control patients.

We now elaborate on the variables included in $\boldsymbol{M}^v$. 

In each visit setting, we observe the following manifest variables:
\begin{enumerate}
    \item $M_1^v(t)$: ICD-10 codes for homelessness and housing instability
    \item $M_2^v(t)$: Document-level classifications of clinical notes for housing status and risk
    \item $\boldsymbol{M}_3^v(t)=[M_{3,1}^v(t), M_{3,2}^v(t), \ldots]$: Counts of spans extracted from clinical notes that indicate homelessness, housing instability, or stable housing
\end{enumerate}

The specific form of the manifest variables depends on the visit setting $v$. In homeless and HUD-VASH visits, specific ICD-10 codes are used to indicate particular housing situations:

\begin{align*}
    M_1^v(t) = \begin{cases}
        0 & \text{No relevant ICD-10 code} \\
        1 & \text{Risk of homelessness} \\
        2 & \text{Unspecified homelessness} \\
        3 & \text{Sheltered} \\
        4 & \text{Unsheltered} \\
    \end{cases} \quad \text{for } v \in \{1, 2\}
\end{align*}

In clinical visits, ICD-10 codes have less granularity and are used primarily to indicate homelessness in general:
\begin{align*}
    M_1^v(t) = \begin{cases}
        0 & \text{No relevant ICD-10 code} \\
        1 & \text{Any code for homelessness or housing instability} \\
    \end{cases} \quad \text{for } v=3 
\end{align*}

In all three settings, document-level classifications are used to construct a three-level categorical variable that indicates whether the notes from that visit suggest the patient is stably housed, at risk of homelessness, or homeless:

\begin{align*}
    M_2^v(t) = \begin{cases}
        0 & \text{No relevant notes} \\
        1 & \text{Majority of notes classified as ``Stably housed''} \\
        2 & \text{Majority of notes classified as ``Unstably housed''} \\
    \end{cases} \quad \text{for } v \in \{1, 2, 3\}
\end{align*}

Finally, we extract counts of spans from clinical notes that indicate homelessness, housing instability, or stable housing. These span counts are expected to be informative about housing status and to differ in their measurement properties across visit settings. 

\paragraph{Parameterization} We parameterize the measurement model using a set of parameters $\boldsymbol{\theta}$ that include the manifest variable parameters $p_k(h, v)$ for each manifest variable $M_k$ and visit setting $v$, as well as the transition probabilities $\pi(h' \mid h, a, x)$ for each baseline state $h$, treatment group $a$, and program enrollment $x$. 



\paragraph{Baseline calibration.} At $t=0$, housing status $H(0) \in \{0, 1, 2\}$ is known for all patients from the Call Center assessment (recall $P(H(0)=3)=0$), and $X(0)=0$ by the eligibility criteria. Let $M_1$ denote a binary ICD-10 code for homelessness observed during healthcare visits. Define the baseline manifest variable parameter:
\begin{align*}
    p(h) = E[M_1 \mid V(0)=1,\, H(0)=h,\, X(0)=0,\, \boldsymbol{L}]
\end{align*}
which is estimable for $h \in \{0, 1, 2\}$.

Using the observed data, we can estimate parameters for the visit and measurement process:
\begin{align*}
    \begin{split}
    P[V(t)=v \mid H(0)=h, A=a, X(t)=x, \boldsymbol{L}] \\
    E[M_k^v(t) \mid V(t)=v, H(0)=h, A=a, X(t)=x, \boldsymbol{L}]
    \end{split} \text{ for } t \in [0, \min(30, \tau)]
\end{align*}

\paragraph{Naive prediction at follow-up.} For each patient, define the naive predicted probability of $M_1$ at follow-up time $t$ under the assumption that housing status has not changed since baseline:
\begin{align*}
    \widehat{p}_0(h, t) = p(H(0))
\end{align*}

\paragraph{Bias of the naive estimator.} Average across patients within treatment group $A=a$ and program enrollment $X(\tau)=x$, and consider the difference between the observed $M_1$ rate at follow-up and the naive prediction, conditioning on baseline housing state and visit type:
\begin{multline*}
    \Delta(h, a, x) = E\big[M_1(\tau) - p(H(0)) \mid \\
    A=a,\, V(0)=1,\, X(\tau)=x,\, H(0)=h\big]
\end{multline*}

Expand the observed $M_1$ rate by the law of total expectation over the unobserved housing state at $\tau$. Let $\mathcal{H} = \{0, 1, 2, 3\}$ denote the set of possible housing states. Then:
\begin{multline*}
    E[M_1(\tau) \mid V(\tau)=1,\, H(0)=h,\, A=a,\, X(\tau)=x,\, \boldsymbol{L}] \\
    = \sum_{h' \in \mathcal{H}} p(h') \; \pi(h' \mid h,\, a,\, x)
\end{multline*}
where we assume the measurement model is time-invariant:
\begin{align*}
    E[M_1(\tau) \mid V(\tau)=1,\, H(\tau)=h',\, \boldsymbol{L}] = p(h')
\end{align*}
and $\pi(h' \mid h,\, a,\, x) = P(H(\tau)=h' \mid H(0)=h,\, A=a,\, X(\tau)=x,\, \boldsymbol{L})$ denotes the housing transition probability. Note that $p(3)$ is not estimable from baseline data since $P(H(0)=3)=0$.

Substituting, and noting that the transition probabilities sum to one:
\begin{align*}
    \Delta(h, a, x) &= \sum_{h' \in \mathcal{H}} p(h') \; \pi(h' \mid h,\, a,\, x) - p(h) \\
    &= \sum_{h' \in \mathcal{H}} \big[p(h') - p(h)\big] \; \pi(h' \mid h,\, a,\, x)
\end{align*}
The bias equals the transition-probability-weighted difference in manifest variable parameters across housing states.

\paragraph{Path to corrected estimates.} Since $|\mathcal{H}|=4$, the transition probabilities $\pi(h' \mid h, a, x)$ have three free parameters (the fourth is determined by the sum-to-one constraint). Additionally, $p(3)$ is not observed at baseline and must be treated as a further unknown. A single manifest variable $M_1$ provides one equation per baseline state, which is not sufficient. However, with $K \geq 3$ manifest variables $M_1, \ldots, M_K$, we obtain a system of $K$ equations:
\begin{multline*}
    \Delta_k(h, a, x) = \sum_{h' \in \mathcal{H}} \big[p_k(h') - p_k(h)\big] \; \pi(h' \mid h,\, a,\, x), \\
    k=1,\ldots,K
\end{multline*}
where $p_k(h) = E[M_k \mid V=1,\, H=h,\, X=0,\, \boldsymbol{L}]$. With $K \geq 3$ manifest variables whose measurement parameters are sufficiently distinct across housing states, the transition probabilities and the unobserved parameters $p_k(3)$ are identifiable and can be used to compute $E[Y(a) \mid S(a)=s]$ within each stratum.

\subsubsection{Step 5: Generalizing the measurement model}

The simplified model above assumes that the manifest variable parameter $p(h)$ depends only on housing status. In practice, several features of the data-generating process complicate this picture.

\paragraph{Visit-specific measurement.} The relationship between manifest variables and housing status differs across clinical settings. A homelessness ICD code recorded during a visit with homeless clinic staff ($V=1$) carries different information than the same code recorded during a visit with other providers ($V=3$). Accordingly, we generalize the measurement model to depend on visit type:
\begin{align*}
    p_k(h, v) = E[M_k \mid V(t)=v,\, H(t)=h,\, X(t),\, \boldsymbol{L}]
\end{align*}

\paragraph{Visit processes.} Let $N_v(t) = \sum_{u=0}^{t} \mathbf{1}[V(u)=v]$ denote the cumulative count of visits of type $v$ through time $t$, and let $N(t) = \sum_{v=1}^{3} N_v(t)$ denote total visits. The probability of visiting a given setting depends on housing status, treatment, and program enrollment:
\begin{align*}
    \lambda_v(t) = P(V(t)=v \mid H(t),\, A,\, X(t),\, \boldsymbol{L})
\end{align*}
Some patients may have no visits during the follow-up window ($N(\tau)=0$), and this is more likely among patients who initiate SSVF and among those who become stably housed ($H(t)=3$). This creates informative visit censoring: the subset of patients with observed manifest variables at follow-up is not representative of all patients.

\paragraph{Time-varying measurement error.} The baseline calibration data are collected in the context of a Call Center contact, which triggers visits that are specifically focused on housing assessment. The measurement model at baseline may therefore differ from the measurement model at later times. Define the full time-dependent measurement parameter:
\begin{align*}
    p_k(h, v, t) = E[M_k \mid V(t)=v,\, H(t)=h,\, X(t),\, \boldsymbol{L}]
\end{align*}
The simplified model in Step 4 assumed $p_k(h, v, t) = p_k(h)$ for all $v$ and $t$. More generally, the measurement model at baseline ($t=0$) may not equal the measurement model at follow-up ($t=\tau$), so that $p_k(h, v, 0) \neq p_k(h, v, \tau)$.

\paragraph{Differential error across programs.} Patients enrolled in different homelessness programs have different documentation patterns and visit processes. A patient enrolled in HUD-VASH ($X(t) \in \{1, 2\}$) has regular contacts with HUD-VASH staff that generate visit-specific manifest variables, while a patient not enrolled in any program ($X(t)=0$) may only be documented during general healthcare visits. The measurement model should therefore condition on program enrollment:
\begin{align*}
    p_k(h, v, t, x) = E[M_k \mid V(t)=v,\, H(t)=h,\, X(t)=x,\, \boldsymbol{L}]
\end{align*}

\paragraph{Summary.} The full measurement model $p_k(h, v, t, x)$ and visit intensity $\lambda_v(t)$ together determine what manifest variable data are observed at follow-up and how they relate to the latent housing state. The bias expression from Step 4 generalizes to:
\begin{multline*}
    \Delta_k(h, v, a, x) = E\big[M_k(\tau) - p_k(H(0), v, 0, 0) \mid \\
    V(\tau)=v,\, A=a,\, X(\tau)=x,\, H(0)=h\big] \\
    = \sum_{h' \in \mathcal{H}} p_k(h', v, \tau, x) \; \pi(h' \mid h,\, a,\, x) - p_k(h, v, 0, 0)
\end{multline*}
which now depends on both the housing transition probabilities $\pi$ and the change in measurement parameters between baseline and follow-up. Identifying the transition probabilities from these equations requires assumptions that constrain the relationship between the baseline and follow-up measurement models. We develop these assumptions next.

\subsection{Measurement model assumptions for $S=5$}

We begin by developing assumptions for stratum $S=5$ (never enrolled in any VA homeless program), restricting attention to visits with other healthcare providers ($V=3$). Since these patients have no VA homeless program contacts, their documentation patterns are driven by general clinical encounters.

\paragraph{SSVF exit status.} For patients who initiated SSVF ($A=1$), define the SSVF exit status $Z \in \{0, 1, 2\}$:
\begin{itemize}
    \item $Z=0$: Still enrolled in SSVF at $\tau$
    \item $Z=1$: Exited SSVF to unstable housing
    \item $Z=2$: Exited SSVF to stable housing
\end{itemize}

\paragraph{Manifest variables.} Let $M_1(t) \in \{0, 1\}$ denote a binary ICD-10 code for homelessness recorded at time $t$, and let $M_2(t) \in \{0, 1, 2, \ldots\}$ denote the count of homeless-related spans in a patient's clinical note at time $t$.

\paragraph{Conditional independence.} For $t > 0$, manifest variables are independent of treatment, SSVF exit status, and baseline housing given current housing status and baseline covariates:
\begin{align*}
    \boldsymbol{M}(t) \perp A,\, Z,\, H(0) \mid H(t),\, \boldsymbol{L},\, V=3,\, S=5
\end{align*}
This assumes that, conditional on current housing status, the way housing information is recorded during general healthcare visits does not depend on whether the patient was assigned to SSVF, how they exited SSVF, or their housing status at the time of the Call Center contact.

\paragraph{Ordered probit structural model for $H(\tau)$.} We model housing status at follow-up using an ordered probit. Define a latent continuous housing index:
\begin{align*}
    H^*(\tau) = \beta_0 + \beta_1\, H(0) + \beta_2\, A + \boldsymbol{\beta}_3'\, \boldsymbol{L} + \varepsilon, \quad \varepsilon \sim N(0, 1)
\end{align*}
with cutpoints $-\infty = c_0 < c_1 < c_2 < c_3 < c_4 = \infty$ such that:
\begin{align*}
    H(\tau) = j \quad \text{if} \quad c_j < H^*(\tau) \leq c_{j+1}, \quad j = 0, 1, 2, 3
\end{align*}
This yields transition probabilities:
\begin{align*}
    \pi(j \mid h, a) &= \Phi(c_{j+1} - \mu(h, a)) - \Phi(c_j - \mu(h, a))
\end{align*}
where $\mu(h, a) = \beta_0 + \beta_1 h + \beta_2 a + \boldsymbol{\beta}_3' \boldsymbol{L}$ and $\Phi$ is the standard normal CDF. The ordered probit parameterizes the full $4 \times 4$ transition matrix through a small number of regression coefficients ($\beta_0, \beta_1, \beta_2, \boldsymbol{\beta}_3$) and three free cutpoints ($c_1, c_2, c_3$), rather than requiring separate transition probabilities for each $(h, a)$ combination.

\paragraph{Measurement model for $M_1$.} We model the probability of a homelessness ICD code using a logistic regression. Let $t \in \{0, 1\}$ index the time period, with $t=0$ for baseline and $t=1$ for follow-up. Then:
\begin{multline*}
    \text{logit}\, P(M_1(t) = 1 \mid H(t),\, \boldsymbol{L},\, V=3,\, S=5) \\
    = \gamma_{10} + \gamma_{11}\, H(t) + \gamma_{12}\, t + \gamma_{13}\, H(t) \cdot t + \boldsymbol{\gamma}_{14}'\, \boldsymbol{L} + \boldsymbol{\gamma}_{15}'\, \boldsymbol{L} \cdot H(t)
\end{multline*}
The terms $\gamma_{12}$ and $\gamma_{13}$ allow the intercept and slope in $H(t)$ to shift between baseline and follow-up, capturing time-varying sensitivity and specificity. The terms $\boldsymbol{\gamma}_{14}'\boldsymbol{L}$ and $\boldsymbol{\gamma}_{15}'\boldsymbol{L} \cdot H(t)$ allow baseline covariates to modify the measurement model. The absence of $\boldsymbol{L} \cdot t$ interactions encodes the assumption that differential measurement error by baseline covariates is time-invariant.

\paragraph{Measurement model for $M_2$.} We model the count of homeless-related note spans using a log-linear model, conditional on $M_1(t)$:
\begin{multline*}
    \log\, E(M_2(t) \mid H(t),\, M_1(t),\, \boldsymbol{L},\, V=3,\, S=5) \\
    = \gamma_{20} + \gamma_{21}\, H(t) + \gamma_{22}\, t + \gamma_{23}\, H(t) \cdot t \\
    + \boldsymbol{\gamma}_{24}'\, \boldsymbol{L} + \boldsymbol{\gamma}_{25}'\, \boldsymbol{L} \cdot H(t) + \gamma_{26}\, M_1(t) + \gamma_{27}\, M_1(t) \cdot H(t)
\end{multline*}
Conditioning on $M_1(t)$ captures within-visit dependence between the two manifest variables. The interaction $M_1(t) \cdot H(t)$ allows the relationship between the ICD code and note spans to vary by housing status: a patient with a homelessness ICD code and many homeless-related note spans is more informative of actual homelessness than a patient with the ICD code but few such spans.
